% Generated draft: requires mathematical and source review.
\documentclass[11pt]{article}
\usepackage[margin=25mm]{geometry}
\usepackage{fontspec}
\setmainfont{Latin Modern Roman}
\usepackage{amsmath,amssymb,mathtools,unicode-math}
\setmathfont{Latin Modern Math}
\usepackage{enumitem,needspace,xurl,hyperref}
\setlength{\parindent}{0pt}
\setlength{\parskip}{5pt}
\setlength{\emergencystretch}{3em}
\newcommand{\N}{\mathbb{N}}
\newcommand{\Z}{\mathbb{Z}}
\newcommand{\Q}{\mathbb{Q}}
\newcommand{\R}{\mathbb{R}}
\newcommand{\C}{\mathbb{C}}
\newcommand{\F}{\mathbb{F}}
\newcommand{\A}{\mathbb{A}}
\newcommand{\PP}{\mathbb{P}}
\newcommand{\E}{\mathbb{E}}
\newcommand{\eps}{\varepsilon}
\newcommand{\dd}{\,\mathrm{d}}
\begin{document}
\section*{UnsolvedMath ID: 801029}
\textbf{Source problem code:} catalogue entry
\subsection{Definitions and mathematical statement}
For a set \(X\subseteq\N\), write \(X\leq_T Y\) for Turing reducibility and \(X'\) for the Turing jump. The set \(X\) is \(\Delta^0_2\) if \(X\leq_T\varnothing'\), and is low if \(X'\equiv_T\varnothing'\). It is \(1\)-generic if, for every computably enumerable set \(W\subseteq 2^{<\N}\) of finite binary strings, either some initial segment of \(X\) belongs to \(W\), or some initial segment of \(X\) has no extension belonging to \(W\).

The supplied source attributes to [LMNS, Section 3.1] the implication chain
\[
\text{\(1\)-generic \(\Delta^0_2\)}
\Longrightarrow \text{index guessable}
\Longrightarrow \text{computes no maximal tower}
\Longrightarrow \text{low}.
\]
It asks whether the first and second implications are reversible.

To display the requested converses without supplying missing hypotheses, let \(\mathcal U\) denote the source's intended, unspecified domain, and let \(G(x)\), \(I(x)\), and \(T(x)\) stand respectively for its predicates ``\(1\)-generic \(\Delta^0_2\)'', ``index guessable'', and ``computes no maximal tower''. The two questions are whether
\[
(\forall x\in\mathcal U)\,[I(x)\Longrightarrow G(x)]
\qquad\text{and}\qquad
(\forall x\in\mathcal U)\,[T(x)\Longrightarrow I(x)]
\]
hold, separately.

\emph{Unresolved formulation qualifications.} The supplied record does not define index guessability or maximal towers, specify the objects and order underlying the towers, or explain how an oracle computes a tower. It also does not specify \(\mathcal U\), including whether its assertions concern individual sets or Turing degrees. For individual sets, the first predicate has the definition given above. A degree formulation could instead ask whether a degree contains a \(1\)-generic \(\Delta^0_2\) representative; that is a different formulation and is not selected here. Thus the displayed converse questions are a faithful logical schema, pending the missing definitions and domain from [LMNS, Section 3.1], rather than a fully specified mathematical conjecture.
\subsection{Short English statement}
The source asks whether index guessability implies the preceding \(1\)-generic \(\Delta^0_2\) condition, and whether computing no maximal tower implies index guessability. The intended domain and the specialized definitions are absent, so both converse questions require clarification from the cited source.
\subsection{Sources}
\begin{itemize}
\item Local UnsolvedMath registry, permanent ID \(801029\), exact source problem code \texttt{catalogue entry}, title \emph{Open Problems in Mathematical Logic --- Problem 41}. The supplied registry labels the problem open; this label is not an independently verified mathematical or literature assessment.
\item Barmpalias, Gao et al., \emph{Open Problems in Mathematical Logic} (2026), Problem 41, \url{https://arxiv.org/abs/2608.26628}. The registry attributes the problem to Barmpalias, Gao et al. and gives proposed year 2026. Its provenance records access on 2026-10-16 and extraction by \texttt{arxiv\_latex}; these are supplied provenance statements, not a claim of independent access or verification.
\item [LMNS, Section 3.1], cited in the supplied statement as establishing the displayed implications. The record supplies no full bibliographic identification or link for [LMNS], so its definitions, hypotheses, and results cannot be checked from the supplied material alone.
\end{itemize}
\end{document}
